% ---------------------------------------------------------------------
% Assembly reference for the remaining manuscript work: for each item
% still to be written, what theory exists, what numbers to quote, which
% script and figure produce them, and what is still open.
%
% Every number here has been evaluated in this repository; the script
% that produces it is named beside it.
% ---------------------------------------------------------------------

\documentclass[10pt]{article}
\usepackage[a4paper,margin=20mm]{geometry}
\usepackage{amsmath,amssymb,booktabs,array,longtable}
\setlength{\parskip}{0.35em}\setlength{\parindent}{0pt}
\newcommand{\scr}[1]{\texttt{\small #1}}
\title{\vspace{-16mm}\textbf{Assembly reference: what is ready, what is
open}\vspace{-3mm}}
\date{\small companion to \emph{cosh\_methodology},
\emph{access\_tight\_rms\_bound}, \emph{ge\_offset\_shift},
\emph{noise\_model\_notes}}
\begin{document}\maketitle

\section{Item A --- why the moment pair may be averaged}

\emph{Status: complete.} Theory, closed form and measurement all
exist; only the manuscript prose is left.

\paragraph{The general statement.} Write the unmodelled part of the
onset balance as $\Delta_\pm$ and split it,
$\Delta_a = (\Delta_+-\Delta_-)/2$, $\Delta_s = (\Delta_++\Delta_-)/2$.
\emph{Everything that acts as force $\times$ pivot distance changes
sign with the tip direction} and therefore lands in $\Delta_a$, so
\begin{equation*}
  M_{\mathrm{ff}} = \tfrac12 (M_+ + M_-)
  = \tfrac12 (W-f)(l_+-l_-) + S_{\mathrm{off}} + \Delta_s ,
\end{equation*}
in which $\Delta_a$ has cancelled identically. The ground-effect
thrust channel, a common displacement of the contact resultant and
rolling resistance are all instances. Measured:
$|\Delta_a| = 54$ mN\,m and $|\Delta_s| = 37$ mN\,m at the median, the
latter $3.7\%$ of the half-width of the interval $[M_-,M_+]$ the
excitation measures directly ($1000$ mN\,m median, $8.2\times$ the
worst $|\Delta_s|$). \hfill (\emph{cosh\_methodology}, present text)

\paragraph{The closed form for the offset.} With the matched-thrust
gains $(\alpha_f,\alpha_M)$,
\begin{equation*}
  \hat p_{\mathrm{off},s} - p_{\mathrm{off}}
  = -\frac{\alpha_M (p_{\mathrm{off}} + s\,l_s)}{1+\alpha_M}
    + \frac{s\,f l_s (\alpha_M-\alpha_f)}{W(1+\alpha_M)},
\end{equation*}
\begin{equation*}
  \hat p_{\mathrm{off,avg}} - p_{\mathrm{off}}
  = -\frac{\alpha_M\,p_{\mathrm{off}}}{1+\alpha_M}
    - \Bigl[\alpha_M - \frac{f(\alpha_M-\alpha_f)}{W}\Bigr]
      \frac{l_+-l_-}{2(1+\alpha_M)} .
\end{equation*}
The slope in $p_{\mathrm{off}}$ is $-\alpha_M/(1+\alpha_M) = -4.13\%$
--- the moment channel alone; the thrust channel only scales the
arm-asymmetry constant and vanishes at $\alpha_f=\alpha_M$, where the
pair average reduces to Eq.~(120). Over the $\pm 20$ mm design box a
single direction carries $\pm 4.1$--$5.8$ mm while the pair average
stays within $\pm 0.83$ mm.
\hfill (\scr{matlab/offset\_ge\_theory.m}, \emph{ge\_offset\_shift})

\paragraph{The measurement.} Per direction $7.5$ / $6.4$ mm RMS
(pooled $7.0$, worst $12.2$); paired $1.77$ mm RMS (worst $2.96$),
with per-configuration $95\%$ confidence intervals of $0.37$--$1.73$
mm (median $0.66$). Ground effect alone predicts $5.4$ / $5.1 \to
0.48$ mm, and $12$ of the $20$ single-direction estimates fall outside
that bound --- so the pairing removes more than ground effect: it
removes every antisymmetric source at once.
\hfill (\scr{offset\_experimental\_split.py}, \scr{fig\_offset\_experimental})

\paragraph{The truth it is measured against.} Load-cell standard error
$0.003$--$0.64$ mm over the ten case-axis measurements, pooled
$0.30$, median $0.03$; datasheet propagation $u_p = l u_f/W = 0.15$ mm.
Contributes $2.8\%$ of the variance of the paired identification error
and is uncorrelated with it ($r = -0.08$).
\hfill (\scr{loadcell\_uncertainty.py}, \scr{DataSet/loadcell})

\section{Item B --- the calibration in the introduction}

\emph{Status: pipeline settled; prose to write.}

The order is worth stating exactly, because it is easy to read
backwards. \emph{Stage 1:} per-run free nonlinear least squares
(PNLS) over $(C_1,C_2,C)$; the per-dataset \emph{medians} of
$(C_2,K)$ locate a search centre. \emph{Stage 2:} the
ramp-invariance score is minimised on a grid around that centre,
fixing $(C_2,K)$; these are frozen in \scr{pnls\_constants.py}.
\emph{Then} COSH runs with the amplitude pinned at $C_1 = K\dot M$,
the baseline at the median, and the onset swept exhaustively. PNLS
seeds the calibration; it is not part of the estimator.

\paragraph{The robustness sentence to include.} The \scr{cosh\_wide}
variant minimises the same score over the wide box
$[3,8]\times[0.05,0.70]$ \emph{with no PNLS centre} and is visually
indistinguishable from COSH. This is the one line that forecloses
``the result depends on the PNLS seed''.

\section{Item C --- the Gazebo ablation}

\emph{Status: open. This is the only item with no material yet.}

\paragraph{What it should target.} The hardware cannot separate the
static ground effect from a contact-resultant shift: a $19$ mm inboard
shift closes the same threshold gap that the interference model
closes. Simulation can set the two independently, so the ablation is
worth most if aimed there rather than at re-confirming what the
hardware already shows.

\begin{center}\small
\begin{tabular}{@{}llp{0.46\linewidth}@{}}
\toprule
run & configuration & what it establishes\\
\midrule
A & GE off, lever exact &
  the pipeline is unbiased: identified offset $=$ truth\\
B & GE on, lever exact &
  the gap appears, and its size is testable against the
  interference model's $\approx 64$ mN\,m residual level\\
C & GE off, lever shifted $19$ mm &
  the same gap reappears --- the degeneracy demonstrated
  directly rather than argued\\
\addlinespace
D & slow ramp, $u>1$ &
  optional: exercises the admissibility floor of (9a),
  $\dot M > \bar\rho C_2$\\
\bottomrule
\end{tabular}
\end{center}

A, B and C together let the manuscript say: \emph{the two mechanisms
are individually sufficient and jointly unidentifiable at this
clearance}, which is stronger than the present argument and costs no
hardware.

\section{Item D --- figures and tables already available}

\begin{center}\small
\begin{tabular}{@{}lll@{}}
\toprule
figure & script & carries\\
\midrule
\scr{fig\_estimator\_err} & \scr{estimator\_error\_figure.py} &
  offset error, six estimators, $95\%$ CI\\
\scr{fig\_mcrit\_static} & \scr{mcrit\_static\_figure.py} &
  thresholds vs three physics levels\\
\scr{fig\_mcrit\_report} & \scr{mcrit\_report\_figure.py} &
  parity plot $+$ physics ladder\\
\scr{fig\_offset\_experimental} & \scr{offset\_experimental\_split.py} &
  per-direction vs paired offsets\\
\scr{fig\_kernel\_free} & \scr{kernel\_free\_bound.py} &
  the residual cap on all $140$ runs\\
\scr{fig\_sg\_bars} & \scr{sg\_bound\_bars.py} &
  the cap with the declared noise constant\\
\scr{fig\_design\_range} & \scr{design\_range\_chart.py} &
  the bound over the design box\\
\scr{fig\_noise\_spectrum} & \scr{noise\_spectrum\_figure.py} &
  the spectral construction behind (19)\\
\bottomrule
\end{tabular}
\end{center}

Tables come from \scr{mcrit\_prediction.csv} (per case/axis/direction)
and \scr{mcrit\_per\_run.csv} (per run), both written by
\scr{mcrit\_prediction.py}.

\section{Numbers to quote, in one place}

\begin{center}\small
\begin{tabular}{@{}lll@{}}
\toprule
quantity & value & source\\
\midrule
static residual, no GE $\to$ interference &
  $191 \to 64$ mN\,m (per-run median) & \scr{mcrit\_report\_figure.py}\\
\quad of which the thrust matching &
  $143 \to 64$ mN\,m & \scr{mcrit\_matched\_thrust.py}\\
one-sided over-prediction &
  $20/20 \to 12/20$ & \scr{mcrit\_prediction.py}\\
share of the gap explained & $81\%$ &
  \emph{cosh\_methodology}\\
\addlinespace
offset error, single $\to$ paired &
  $7.0 \to 1.77$ mm RMS & \scr{offset\_experimental\_split.py}\\
load-cell truth uncertainty &
  $0.30$ mm pooled ($0.64$ worst) & \scr{loadcell\_uncertainty.py}\\
\addlinespace
residual cap coverage & $140/140$, used $0.36$--$0.41$ &
  \scr{kernel\_free\_bound.py}\\
noise constant $N_n$ & $2.31^\circ$/s $= 3N_{\mathrm{med}}$ &
  \scr{sg\_bound\_bars.py}\\
realised $\delta e_\omega$ & $0.16$--$0.17^\circ$/s &
  \scr{kernel\_free\_bound.py}\\
$\bar\rho$ at the $5^\circ$ box & $1.8$--$3.3$ mN\,m &
  \scr{design\_range\_chart.py}\\
GE share of $\bar\rho$ & $2$--$10\%$; $6\times$ inflation $\to$
  cap $+0.1\%$ & \scr{ge\_channel\_share.py}\\
\bottomrule
\end{tabular}
\end{center}

\section{Two cautions carried forward}

\emph{The dynamic ground effect.} The attitude dependence is not
resolvable here: the inversion's sensitivity floor is
$\approx 90$ mN\,m --- the static repeatability --- against a model
variation of $-2.5$ to $-0.1$ mN\,m/deg over the $5^\circ$ cap. The
bound route is not a fallback but the tighter statement: unmodelled
forcing that the family cannot absorb shows in the residual, capped at
$140/140$; what the family does absorb shifts $M_{\mathrm{crit}}$ by at
most $\bar\rho = 1.8$--$3.3$ mN\,m through (21). Both branches must be
stated together --- the first alone invites ``the fit absorbed it''.

\emph{Attribution language.} The interference model \emph{accounts
for} the gap; it does not \emph{identify} ground effect, because a
$19$ mm contact-resultant shift closes the same gap. Keep
``the size of the near-ground correction, not its attribution''
wherever the static result is stated.

\end{document}
